% !PRINTED 245
\weberpage{245}
%<<CONT
worin $a_{h,k}$ ganze Zahlen sind, die den Bedingungen
\[
a_{h,k} = a_{k,h}, \quad a_{p,p} = 0
\]
genügen.

\begin{center}
\textbf{§ 70.\quad \gesperrt{Transformationsgleichungen erster Stufe.}}
\end{center}

Die Invariantengleichung ist von großer theoretischer Wichtigkeit teils
wegen ihrer allgemeinen Gültigkeit (auch für gerade $n$), teils wegen der
Leichtigkeit, mit der die linearen Transformationen auf $j(\omega)$
angewandt werden können. Die wirkliche Berechnung dieser Gleichung
aber zeigt sich so kompliziert, und die Zahlenkoeffizienten sind so
groß, daß die Berechnung bis jetzt nur in dem einfachsten Falle $p=2$
durchgeführt ist. Dagegen kann man, indem man andere Modulfunktionen
benutzt, weit einfachere Transformationsgleichungen erhalten.

Über das hierbei anzuwendende Prinzip bemerken wir folgendes:

Wenn irgend ein System von $\nu$ Funktionen von $\omega$ vorliegt,
entsprechend den $\nu$ Systemen von Transformationszahlen $a$, $c$,
$\partial$, die wir mit
\weq{1}{\Phi_{a,c,\partial}}
bezeichnen wollen und die \gesperrt{isomorph} mit den Funktionen
\weq{2}{j\left(\frac{c+\partial\omega}{a}\right)}
durch die Substitutionen
\weq{3}{\left(\omega, -\frac{1}{\omega}\right), \quad (\omega, \omega+1)}
untereinander permutiert werden, so ist (1) \gesperrt{rational} durch (2) und
durch $j(\omega)$ ausdrückbar, denn die Funktion

\weq{4}{F_n[x,j(\omega)] \sum^{a,c,\partial} \frac{\Phi_{a,c,\partial}}{x-j\left(\frac{c+\partial\omega}{a}\right)} = \Psi[x,j(\omega)]}

bleibt durch die Substitutionen (3) ungeändert, und ist daher (für ein
unbestimmtes $x$) eine \gesperrt{rationale Funktion} von $j(\omega)$ und überdies
eine ganze rationale Funktion von $x$, höchstens vom Grade $\nu-1$;
$F_n[x,j(\omega)]$ hat dieselbe Bedeutung, wie in (17) des vorigen
Paragraphen, und $\nu=\psi(n)$ ist der Grad dieser Funktion in bezug
auf $x$. Aus (4) aber erhält man, indem man
\[
x = j\left(\frac{c+\partial\omega}{a}\right)
\]
setzt:
%CONT>>
